\subsection{Clustering and Mixture Models}

Unsupervised learning searches for structure in inputs without direct labels.  Clustering is the simplest form of latent assignment.

\begin{definition} \label{def:sl-kmeans-objective} \hlt{$k$-Means Clustering}\\
Given data $x_1,\ldots,x_n\in\R^d$, $k$-means chooses centers $\mu_1,\ldots,\mu_k$ and assignments $c_i\in\{1,\ldots,k\}$ to minimize
\[
\sum_{i=1}^n\norm{x_i-\mu_{c_i}}_2^2.
\]
\end{definition}

\begin{definition} \label{def:sl-finite-mixture-model} \hlt{Finite Mixture Model}\\
A finite mixture has density
\[
p(x)=\sum_{z=1}^k\pi_z p_z(x;\theta_z),
\qquad
\pi_z\geq0,\quad \sum_{z=1}^k\pi_z=1,
\]
where the component label $Z$ is latent.
\end{definition}

\begin{method} \label{method:sl-em-algorithm} \hlt{Expectation-Maximization}\\
The EM algorithm alternates between computing posterior responsibilities for latent variables under current parameters and maximizing the expected complete-data log-likelihood.  It is a coordinate ascent method for likelihood lower bounds.
\end{method}

\begin{remark} \label{rem:sl-clustering-identifiability} \hlt{Identifiability}\\
Mixture components are identifiable only up to label permutations.  This is a benign symmetry for density estimation but important when interpreting clusters.
\end{remark}
